%##########################################################################
% CHAPTER 3: METHODOLOGY
%##########################################################################
\chapter{Methodology}

\section{Chapter Overview}

This chapter lays out the research design, walks through each of the five system modules in detail, describes the testbed game, and specifies the evaluation plan. The methodology follows a mixed-methods approach \citep{creswell2017mixed}, combining a controlled between-subjects experiment with qualitative interview analysis to capture both the measurable and the experiential dimensions of narrative personalization.

\section{Research Design Overview}

The research adopts a mixed-methods design that pairs a controlled between-subjects experiment with a qualitative interview study. The experimental component is designed for causal inference: if a significant difference in player experience is observed between conditions, the controlled design allows it to be attributed to the adaptive narrative system rather than to confounding factors. The qualitative component captures something the numbers cannot --- the texture of how players actually experience personalization, what they notice, what they value, and what fails to land.

A between-subjects design was chosen over a within-subjects alternative to eliminate order effects and carry-over contamination. The concern is straightforward: a player who first experiences adaptive narrative and then switches to static narrative (or vice versa) would inevitably bring expectations, comparisons, and memory from the first condition into the second. The experience of the second condition would be coloured by the first, and any difference between conditions would be confounded with the order in which they were encountered. With a between-subjects design, each participant experiences exactly one condition, cleanly eliminating this contamination at the cost of requiring a larger sample and careful matching to ensure the two groups are comparable at baseline.

The experimental condition (Experimental) deploys the full five-module AI system described in this chapter. The control condition (Control) replaces the Narrative Generation Service and the Adaptation Engine with a \texttt{StaticNarrativeContent} ScriptableObject --- a pre-authored collection of narrative fragments delivered in response to the same game triggers, in the same visual format, but with no dependence on player state. Everything else --- the game environment, the objectives, the NPC placement, the session duration, the UI layout --- is identical across conditions. This isolation is critical: if the Experimental group differs from the Control group on the dependent measures, the difference can be attributed to the adaptive narrative system itself, not to differences in game content, visual presentation, or session structure.

\section{System Architecture}

The system is organized as five modules in a closed-loop pipeline. Figure~\ref{fig:architecture} describes the architecture and the data flows between components.

\begin{figure}[H]
\centering
\resizebox{\textwidth}{!}{%
\begin{tikzpicture}[
    node distance=1.2cm and 1.8cm,
    >={Stealth[length=3mm]},
    module/.style={
        rectangle, rounded corners=4pt, draw=black!70, line width=0.8pt,
        fill=#1, minimum width=3.2cm, minimum height=1.6cm,
        text width=3cm, align=center, font=\small
    },
    module/.default={blue!8},
    component/.style={
        rectangle, rounded corners=2pt, draw=black!50, line width=0.5pt,
        fill=white, minimum width=2.4cm, minimum height=0.55cm,
        align=center, font=\scriptsize\ttfamily
    },
    data/.style={
        font=\scriptsize\itshape, text=black!70
    },
    arrow/.style={
        ->, line width=0.9pt, black!70
    },
    groupbox/.style={
        rectangle, rounded corners=6pt, draw=black!40, dashed,
        line width=0.6pt, inner sep=10pt
    }
]

% ── Module 1: Unity Client (Telemetry) ──
\node[module=green!10] (telemetry) {
    \textbf{Module 1}\\[2pt]
    Telemetry\\Collection
};
\node[component, below=0.3cm of telemetry] (pdl) {PlayerDataLogger};

% ── Module 2: Player Modelling ──
\node[module=blue!10, right=2.5cm of telemetry] (playermod) {
    \textbf{Module 2}\\[2pt]
    Player\\Modelling
};
\node[component, below=0.15cm of playermod, yshift=-0.1cm] (fe) {FeatureExtractor};
\node[component, below=0.05cm of fe] (hp) {HexadProfiler};
\node[component, below=0.05cm of hp] (fc) {FlowClassifier};

% ── Module 3: Adaptation Engine ──
\node[module=orange!12, right=2.5cm of playermod] (adaptation) {
    \textbf{Module 3}\\[2pt]
    Adaptation\\Engine
};
\node[component, below=0.3cm of adaptation] (ppo) {NarrativeAgent (PPO)};

% ── Module 4: Narrative Generation ──
\node[module=purple!10, below=3.2cm of adaptation] (narrative) {
    \textbf{Module 4}\\[2pt]
    Narrative\\Generation
};
\node[component, below=0.15cm of narrative, yshift=-0.1cm] (rag) {LoreRetriever (RAG)};
\node[component, below=0.05cm of rag] (mem) {EpisodicMemory};
\node[component, below=0.05cm of mem] (pb) {PromptBuilder};
\node[component, below=0.05cm of pb] (ollama) {OllamaClient};

% ── Module 5: Content Injection ──
\node[module=green!10, below=3.2cm of telemetry] (injection) {
    \textbf{Module 5}\\[2pt]
    Content\\Injection
};
\node[component, below=0.15cm of injection, yshift=-0.1cm] (nm) {NarrativeManager};
\node[component, below=0.05cm of nm] (ci) {ContentInjector};

% ── Player node ──
\node[module=gray!15, left=2.0cm of telemetry, yshift=-1.6cm] (player) {
    \textbf{Player}\\[2pt]
    {\scriptsize Gameplay}\\
    {\scriptsize Behaviour}
};

% ── Backend group box ──
\begin{scope}[on background layer]
    \node[groupbox, fill=yellow!4,
          fit=(playermod)(fe)(hp)(fc)(adaptation)(ppo)(narrative)(rag)(mem)(pb)(ollama),
          label={[font=\small\bfseries, text=black!60]above:FastAPI Backend}
    ] (backend) {};
\end{scope}

% ── Unity group box ──
\begin{scope}[on background layer]
    \node[groupbox, fill=green!4,
          fit=(telemetry)(pdl)(injection)(nm)(ci),
          label={[font=\small\bfseries, text=black!60]above:Unity Client}
    ] (unity) {};
\end{scope}

% ── Arrows ──
% Player -> Telemetry
\draw[arrow] (player.north east) -- node[data, above, sloped] {gameplay} (telemetry.west);

% Telemetry -> Player Modelling
\draw[arrow] (telemetry.east) -- node[data, above] {TelemetryBatch} node[data, below] {\scriptsize HTTP POST / 5s} (playermod.west);

% Player Modelling -> Adaptation
\draw[arrow] (playermod.east) -- node[data, above] {PlayerModel} (adaptation.west);

% Adaptation -> Narrative Generation
\draw[arrow] (adaptation.south) -- node[data, right, text width=2cm] {NarrativeAction} (narrative.north);

% Narrative Generation -> Content Injection
\draw[arrow] (narrative.west) -- node[data, above] {NarrativeContent} node[data, below] {\scriptsize JSON response} (injection.east);

% Content Injection -> Player
\draw[arrow] (injection.west) -| node[data, below, near start] {dialogue / UI} (player.south);

% ── Lore corpus annotation ──
\node[component, fill=purple!5, right=0.8cm of rag, minimum width=2cm] (lore) {Lore Corpus};
\draw[->, dashed, black!50] (lore) -- (rag);

% ── Ollama annotation ──
\node[component, fill=orange!5, right=0.8cm of ollama, minimum width=2cm] (phi) {Phi-3.5 Mini};
\draw[->, dashed, black!50] (phi) -- (ollama);

\end{tikzpicture}%
}
\caption{Closed-loop system architecture. Five modules form a pipeline from player behaviour through telemetry collection, player modelling, PPO-based adaptation, RAG-grounded narrative generation, and content injection back to the player. Dashed boxes indicate platform boundaries.}
\label{fig:architecture}
\end{figure}

The pipeline operates as a continuous loop, cycling once per telemetry window. The Unity \texttt{PlayerDataLogger} aggregates behavioural signals over five-second windows and transmits a \texttt{TelemetryBatch} to the FastAPI backend via HTTP POST (with WebSocket support available for lower-latency streaming when network conditions permit). On the backend, the Player Modelling Service processes each incoming batch through three sequential stages: the \texttt{FeatureExtractor} computes eleven normalised features from the raw telemetry, the \texttt{HexadProfiler} derives the six-dimensional continuous Hexad profile from the cumulative session history, and the \texttt{FlowClassifier} assigns a Flow state based on the most recent data. The assembled \texttt{PlayerModel} is then passed to the \texttt{NarrativeAgent}, which selects a \texttt{NarrativeAction} --- using the cold-start heuristic during the first 120~seconds, and the trained PPO policy thereafter. The Narrative Generation Service takes the selected action and uses it to formulate a query for the \texttt{LoreRetriever} (RAG), consults the \texttt{EpisodicMemory} for relevant prior narrative events, assembles the full two-message prompt through the \texttt{PromptBuilder}, and dispatches a generation request to the locally hosted Ollama LLM. The resulting \texttt{NarrativeContent} --- a JSON payload containing the generated text, content type, and speaker attribution --- is returned to the Unity client and injected at the appropriate in-game display point by the \texttt{ContentInjector}.

\section{Player Modelling Module}

The player modelling module is responsible for answering two questions about the player at every point in the session: \textit{what kind of player are they?} (the Hexad profile) and \textit{how are they doing right now?} (the Flow state). Everything downstream --- the action the PPO agent selects, the tone the LLM adopts, the content the player ultimately sees --- depends on the accuracy and responsiveness of these two answers.

\subsection{Telemetry Collection}

The telemetry system captures 23 behavioural fields per five-second window, providing the raw signal from which both the Hexad profile and the Flow state are derived. Table~\ref{tab:telemetry} lists the complete schema.

\begin{longtable}{@{}p{4cm} p{1.5cm} p{7.5cm}@{}}
\caption{TelemetryBatch schema (23 fields).} \label{tab:telemetry} \\
\toprule
\textbf{Field} & \textbf{Type} & \textbf{Description} \\
\midrule
\endfirsthead
\toprule
\textbf{Field} & \textbf{Type} & \textbf{Description} \\
\midrule
\endhead
player\_id & str & Unique player identifier \\
session\_id & str & Unique session identifier \\
window\_start & float & Window start time (Unix timestamp) \\
window\_end & float & Window end time (Unix timestamp) \\
kills & int & Enemy kills in window \\
deaths & int & Player deaths in window \\
damage\_taken & float & Total damage received \\
damage\_dealt & float & Total damage inflicted \\
abilities\_used & int & Ability activations \\
abilities\_hit & int & Abilities that connected \\
objectives\_completed & int & Objectives completed \\
objectives\_attempted & int & Objectives attempted \\
tiles\_explored & int & Map tiles visited \\
total\_tiles & int & Total traversable tiles (default 100) \\
npc\_interactions\_voluntary & int & Voluntary (unsolicited) NPC interactions \\
dialogue\_lines\_shown & int & Dialogue lines displayed \\
dialogue\_lines\_skipped & int & Dialogue lines skipped by player \\
lore\_items\_read & int & Lore collectibles read \\
lore\_items\_found & int & Lore collectibles discovered \\
idle\_seconds & float & Seconds with no player input \\
session\_elapsed & float & Total elapsed session time (seconds) \\
current\_location & str & Scene/zone identifier \\
current\_objective & str & Active quest objective identifier \\
\bottomrule
\end{longtable}

The five-second window granularity represents a deliberate trade-off between responsiveness and noise sensitivity. A shorter window would allow faster adaptation cycles but would amplify transient spikes --- a single unlucky death or a momentary pause to answer a phone call would distort the player model. A longer window would produce smoother signals but at the cost of delayed response: if a player tips from Flow into anxiety, the system should notice within seconds, not minutes. Five seconds proved to be a practical sweet spot during development. In practice, the system further smooths the signal by accumulating a sliding window of recent batches (typically six, spanning approximately 30~seconds) before computing the feature vector, allowing multi-window averages to absorb transient fluctuations.

\subsection{Feature Extraction (11 Features)}

The \texttt{FeatureExtractor} aggregates one or more \texttt{TelemetryBatch} records into eleven normalized scalar features, each clipped to [0, 1] (or [0, 3] for the \texttt{challenge\_skill\_ratio}). Table~\ref{tab:features} presents the computation formulae.

\begin{table}[H]
\centering
\caption{Eleven extracted features and computation formulae.}
\label{tab:features}
\small
\begin{tabular}{@{}p{3.5cm} p{6cm} p{3.5cm}@{}}
\toprule
\textbf{Feature} & \textbf{Formula} & \textbf{Notes} \\
\midrule
kd\_ratio & clip(kills / max(deaths, 1) / 5.0, 0, 1) & Norm.\ at 5:1 K/D \\
objective\_rate & clip(obj\_comp / max(obj\_att, 1), 0, 1) & Completion proportion \\
explore\_rate & clip(tiles\_explored / total\_tiles, 0, 1) & Map coverage \\
dialogue\_engage & clip(1 $-$ skipped / max(shown, 1), 0, 1) & 1 = reads all \\
ability\_accuracy & clip(hit / max(used, 1), 0, 1) & Combat precision \\
damage\_taken\_norm & clip(dmg / (elapsed\_min + 1) / 200, 0, 1) & DPM normalized \\
idle\_ratio & clip(idle\_s / elapsed, 0, 1) & Fraction idle \\
lore\_engage\_rate & clip(read / max(found, 1), 0, 1) & Lore depth \\
skill\_score & 0.35 kd + 0.30 obj + 0.20 acc + 0.15(1$-$idle) & Composite skill \\
challenge\_score & 0.50 dmg + 0.30(1$-$obj) + 0.20 death\_ratio & Composite challenge \\
challenge\_skill\_ratio & clip(challenge / max(skill, $10^{-6}$), 0, 3) & Core Flow signal \\
\bottomrule
\end{tabular}
\end{table}

The \texttt{skill\_score} composite weights combat effectiveness (\texttt{kd\_ratio}, \texttt{ability\_accuracy}) most heavily, followed by objective completion and attention (inverse \texttt{idle\_ratio}). The \texttt{challenge\_score} composite weights damage pressure most heavily, followed by objective failure rate and death rate. Dividing challenge by skill yields the \texttt{challenge\_skill\_ratio}, which operationalizes Csikszentmihalyi's (\citeyear{csikszentmihalyi1990flow}) challenge--skill balance axis.

\subsection{Hexad Profile Computation (6 Dimensions)}

The \texttt{HexadProfiler} computes six continuous scores in [0, 1] from accumulated telemetry, each representing the strength of a Hexad player-type signal. Table~\ref{tab:hexad} shows the computation weights.

\begin{table}[H]
\centering
\caption{Hexad dimension computation weights from telemetry signals.}
\label{tab:hexad}
\small
\begin{tabular}{@{}p{2.5cm} p{7cm} p{3.5cm}@{}}
\toprule
\textbf{Dimension} & \textbf{Formula} & \textbf{Interpretation} \\
\midrule
Achiever & 0.70 $\times$ obj\_rate + 0.30 $\times$ norm(kills, 20) & Objective + combat \\
Explorer & 0.50 $\times$ explore\_rate + 0.50 $\times$ lore\_rate & Map + lore coverage \\
Socializer & 0.60 $\times$ norm(vol\_npc, 10) + 0.40 $\times$ dialogue & Seeks NPCs + reads \\
Free Spirit & 0.50 $\times$ explore\_rate + 0.50 $\times$ (1 $-$ obj\_rate) & Explores off-path \\
Disruptor & 0.65 $\times$ norm(kills, 20) + 0.15 $\times$ norm(deaths, 10) + 0.20 $\times$ norm(dmg, 1000) & Aggressive combat \\
Philanthropist & 0.60 $\times$ dialogue\_engage + 0.40 $\times$ lore\_rate & Narrative engagement \\
\bottomrule
\end{tabular}
\end{table}

The mapping from behavioural signals to Hexad dimensions was established through a principled alignment process guided by the original Hexad motivation descriptions \citep{tondello2016hexad}. The logic proceeds dimension by dimension. Achievers seek mastery and progress, so the signal proxies are objective completion rate and combat kill count --- both direct indicators of goal-oriented play. Explorers seek discovery, so map coverage and lore engagement serve as proxies: a player who visits every corner of the map and reads every collectible is behaving as an Explorer regardless of what a questionnaire might say. Socializers seek connection, so voluntary NPC interaction count and dialogue engagement rate are weighted; a player who seeks out NPCs unprompted and reads their dialogue is exhibiting Socializer behaviour. Free Spirits resist imposed structure, so the proxy is exploration \textit{uncorrelated with} objective completion --- wandering off the critical path rather than following it. Disruptors challenge systems, captured by aggressive combat behaviour (high kill counts, high damage dealt, willingness to die repeatedly). Philanthropists give and share, captured by depth of engagement with narrative and lore content --- a proxy for the motivation to understand and connect with the game world rather than merely to conquer it.

All profiles are initialised to 0.5 (the neutral midpoint) in the absence of data, preventing cold-start artefacts from biasing the narrative action selection before the system has observed enough behaviour to form a meaningful profile.

\subsection{Flow State Classification (MDP States)}

The \texttt{FlowClassifier} implements a decision-tree classifier over the \texttt{challenge\_skill\_ratio}, \texttt{idle\_ratio}, and \texttt{dialogue\_engage\_rate} features. Table~\ref{tab:flow_rules} formalizes the rules.

\begin{table}[H]
\centering
\caption{Flow state classification rules (evaluated in priority order).}
\label{tab:flow_rules}
\small
\begin{tabular}{@{}clll@{}}
\toprule
\textbf{Priority} & \textbf{Condition} & \textbf{State} & \textbf{Confidence} \\
\midrule
1 (highest) & idle $>$ 0.40 AND dialogue $<$ 0.30 AND ratio $<$ 0.60 & APATHY & 0.70 \\
2 & ratio $<$ 0.75 & BOREDOM & 0.60 + (0.75 $-$ ratio) $\times$ 0.80 \\
3 & ratio $>$ 1.30 & ANXIETY & 0.60 + (ratio $-$ 1.30) $\times$ 0.40 \\
4 (default) & 0.75 $\leq$ ratio $\leq$ 1.30 & FLOW & max(0.55, 1.0 $-$ $|$ratio $-$ 1.0$|$ $\times$ 1.2) \\
\bottomrule
\end{tabular}
\end{table}

Figure~\ref{fig:flow_tree} visualises the classification logic as a decision tree.

\begin{figure}[H]
\centering
\begin{tikzpicture}[
    scale=0.82, every node/.style={scale=0.82},
    >={Stealth[length=2mm]},
    decision/.style={
        diamond, draw=black!60, line width=0.6pt, fill=gray!6,
        inner sep=2pt, align=center, font=\small,
        minimum width=2.2cm, aspect=2.2
    },
    outcome/.style={
        rectangle, rounded corners=4pt, draw=black!60, line width=0.7pt,
        minimum width=2.2cm, minimum height=0.9cm,
        align=center, font=\small\bfseries, inner sep=5pt
    },
    edgelbl/.style={font=\scriptsize, text=black!70},
]

% ── Decision nodes ──
\node[decision] (d1) {idle $>$ 0.40\\AND dial $<$ 0.30\\AND ratio $<$ 0.60?};
\node[outcome, fill=red!12, right=3.2cm of d1] (apathy) {APATHY};

\node[decision, below=1.6cm of d1] (d2) {ratio $<$ 0.75?};
\node[outcome, fill=yellow!20, right=3.2cm of d2] (boredom) {BOREDOM};

\node[decision, below=1.6cm of d2] (d3) {ratio $>$ 1.30?};
\node[outcome, fill=orange!18, right=3.2cm of d3] (anxiety) {ANXIETY};

\node[outcome, fill=green!15, below=1.6cm of d3] (flow) {FLOW};

% ── Edges ──
\draw[->, line width=0.7pt] (d1) -- node[edgelbl, above] {Yes} (apathy);
\draw[->, line width=0.7pt] (d1) -- node[edgelbl, left] {No} (d2);
\draw[->, line width=0.7pt] (d2) -- node[edgelbl, above] {Yes} (boredom);
\draw[->, line width=0.7pt] (d2) -- node[edgelbl, left] {No} (d3);
\draw[->, line width=0.7pt] (d3) -- node[edgelbl, above] {Yes} (anxiety);
\draw[->, line width=0.7pt] (d3) -- node[edgelbl, left] {No (default)} (flow);

\end{tikzpicture}
\caption{Flow state classification decision tree. Conditions are evaluated in priority order (top to bottom); the first satisfied condition determines the classification.}
\label{fig:flow_tree}
\end{figure}

The priority ordering matters. APATHY is evaluated first because it represents a qualitatively distinct disengagement pattern --- the player has not merely lost interest in the challenge (which would be BOREDOM) but has stopped meaningfully interacting with the game altogether. Without the compound condition check, such a player would be classified as BORED on the challenge--skill ratio alone, and the system would respond with urgency-boosting narrative. But a disengaged player does not need urgency; they need a reason to care. The APATHY classification ensures the system responds with mystery and intrigue rather than pressure.

The confidence scoring for BOREDOM and ANXIETY is proportional to the distance from the FLOW band boundary: a player with a challenge--skill ratio of 0.40 is classified with higher confidence than one at 0.70, because the former is further from the zone of optimal engagement. FLOW confidence is highest when the ratio sits closest to 1.0 --- the theoretical point of perfect challenge--skill balance --- and tapers as the ratio approaches either boundary.

\section{Narrative Generation Module}

The narrative generation module is the part of the system that actually writes the story. Everything upstream --- the telemetry, the feature extraction, the Flow classification, the Hexad profiling, the PPO action selection --- exists to inform this module about \textit{what kind of narrative to produce}. The module itself is responsible for producing it: grounded in the game's lore, consistent with the session's narrative history, and formatted for injection into the game's UI.

\subsection{RAG Pipeline}

The \texttt{LoreRetriever} implements in-memory semantic search using the sentence-transformers all-MiniLM-L6-v2 model, which produces 384-dimensional L2-normalised embeddings. The lore corpus --- comprising the world, character, and quest lore documents stored in the \texttt{lore/} directory --- is pre-indexed at server startup. Each document is chunked by paragraph (double-newline delimiters), and each chunk is embedded into the 384-dimensional space. At retrieval time, the query (constructed from the player's current Flow state, activity, and location) is embedded and cosine similarities are computed against all stored chunk embeddings via a dot product (since all vectors are unit-normalised, cosine similarity reduces to the dot product). The top $k = 3$ chunks by similarity are returned and formatted as bullet points in the LLM user prompt, providing the model with the lore context most relevant to the current generation task.

This design was chosen over a persistent vector database (e.g., ChromaDB, Pinecone) for two reasons that proved decisive during development. First, ChromaDB's Pydantic v1 dependency created compatibility issues with the project's Python 3.14 runtime that could not be cleanly resolved. Second, and more fundamentally, the lore corpus is small --- three files, approximately 30 paragraphs, a few hundred embeddings at most --- which renders the indexing, persistence, and query optimisation machinery of a full vector database unnecessary overhead. In-memory retrieval at this corpus scale completes in sub-millisecond time, well within the latency budget of the narrative generation pipeline.

\subsection{Episodic Session Memory}

The \texttt{EpisodicMemory} class maintains a capped list of \texttt{MemoryEvent} records, each comprising a text description, an importance score in [0, 1], and a timestamp. When the list exceeds the maximum capacity of ten events, it is sorted by importance in descending order and truncated to the top ten. This importance-weighted eviction ensures that the most narratively significant events are retained even as minor events accumulate.

At prompt construction time, the top five events by importance score are selected and formatted as a context block, providing the LLM with a compressed narrative history that enables references to past events and prevents temporal inconsistency.

\subsection{Dynamic Prompt Construction}

The \texttt{PromptBuilder} assembles a two-message prompt (system + user) for the Ollama API. The system prompt encodes the LLM's role (narrative writer for the game world), output length constraints (2--4 sentences for dialogue, 1--2 for descriptions), format requirements (valid JSON only, specific schema), and the emotional tone mapped from the current Flow state:

\begin{itemize}
\item ANXIETY: ``urgent and supportive --- the player is struggling''
\item BOREDOM: ``exciting and surprising --- inject energy''
\item FLOW: ``immersive and atmospheric --- maintain the experience''
\item APATHY: ``intriguing and mysterious --- re-engage the player''
\end{itemize}

The user prompt injects: the player's current Flow state and challenge--skill ratio; the current activity (IN\_COMBAT, EXPLORING, IN\_DIALOGUE, IDLE); session elapsed time; the Hexad profile percentages; the selected NarrativeAction with definitions; the current location and active quest stage; the three retrieved lore chunks; the episodic memory context; and explicit content constraints (no real-world references, no anachronisms, no fourth-wall breaks, character voice descriptions).

\subsection{LLM Integration (Ollama / Phi-3.5 Mini)}

The \texttt{OllamaClient} sends the assembled prompt to a locally hosted Ollama instance using the \texttt{/api/chat} endpoint. The primary model is Phi-3.5 Mini (3.8B parameters; \citealp{microsoft2024phi}), with Llama 3.2 3B as a configured fallback. The generation parameters are: temperature 0.75, top\_p 0.90, maximum tokens 200, JSON format mode enabled. The 8-second HTTP timeout ensures that a slow or stalled generation does not block the main game loop.

When the LLM response is malformed JSON, contains a fallback signal, or the request times out, the \texttt{OllamaClient} returns a pre-authored fallback string corresponding to the requested \texttt{NarrativeAction}. These fallback strings are intentionally minimal --- a single sentence of tonally appropriate narrative --- designed to maintain the illusion of a functioning narrative system without exposing the player to error messages, empty dialogue boxes, or other artefacts that would break immersion. The player sees a brief, contextually appropriate line; they do not see that the LLM failed to produce valid output. This graceful degradation is essential for a system that will be used in a controlled user study, where a single visible error could contaminate a participant's experience and compromise the evaluation.

\section{Adaptation Engine (MDP + PPO)}

The adaptation engine is the decision-making core of the system: given what the player modelling module knows about the player right now, what should the narrative do next? This question is formalised as a reinforcement learning problem, which allows the answer to be \textit{learned} from experience rather than hand-coded by the designer.

\subsection{MDP Formulation}

The narrative adaptation problem is formalised as a Markov Decision Process $(S, A, T, R, \gamma)$ implemented as a Gymnasium environment.

\textbf{State space $S$:} A Box observation space of shape (7,) with all values in [0, 1]. Table~\ref{tab:mdp_state} defines the seven dimensions.

\begin{table}[H]
\centering
\caption{MDP observation vector (7 dimensions).}
\label{tab:mdp_state}
\small
\begin{tabular}{@{}clp{7cm}@{}}
\toprule
\textbf{Index} & \textbf{Signal} & \textbf{Description} \\
\midrule
0 & flow\_score & Encoded Flow state (FLOW=1.0, BOREDOM=0.3, ANXIETY=0.2, APATHY=0.1) \\
1 & csr\_norm & challenge\_skill\_ratio / 3.0 \\
2 & hexad\_explorer & Explorer dimension score \\
3 & hexad\_socializer & Socializer dimension score \\
4 & hexad\_achiever & Achiever dimension score \\
5 & hexad\_disruptor & Disruptor dimension score \\
6 & session\_progress & session\_elapsed / 1800.0 (normalized to 30-min max) \\
\bottomrule
\end{tabular}
\end{table}

\textbf{Action space $A$:} A Discrete(8) action space corresponding to the eight \texttt{NarrativeAction} values. Table~\ref{tab:actions} summarizes the actions.

\begin{table}[H]
\centering
\caption{Eight narrative actions with descriptions and intended effects.}
\label{tab:actions}
\small
\begin{tabular}{@{}cll@{}}
\toprule
\textbf{Idx} & \textbf{Action} & \textbf{Intended Effect} \\
\midrule
0 & LOWER\_STAKES & Reduce narrative tension; provide respite \\
1 & RAISE\_STAKES & Increase narrative urgency and weight \\
2 & ADD\_MYSTERY & Introduce a curiosity hook or unexplained element \\
3 & ADD\_HUMOR & Provide comic relief to reduce tension \\
4 & PROVIDE\_GUIDANCE & Give struggling player directional narrative support \\
5 & INCREASE\_URGENCY & Re-energize a bored or disengaged player \\
6 & LORE\_REWARD & Deliver a lore discovery payoff for curious players \\
7 & NO\_CHANGE & Suppress narrative injection; maintain current experience \\
\bottomrule
\end{tabular}
\end{table}

\textbf{Transition dynamics $T$:} The simulated training environment models stochastic Flow state transitions. When the agent selects a ``helpful'' action for the current state (e.g., ANXIETY $\rightarrow$ \{PROVIDE\_GUIDANCE, LOWER\_STAKES\}), the environment transitions to FLOW with probability 0.75 or remains in the current state with probability 0.25. This stochasticity prevents the policy from memorizing a deterministic mapping and encourages generalization across states.

Figure~\ref{fig:mdp_transitions} visualises the state transitions and the narrative actions that drive them.

\begin{figure}[H]
\centering
\begin{tikzpicture}[
    scale=0.85, every node/.style={scale=0.85},
    >={Stealth[length=2.5mm]},
    state/.style={
        circle, draw=black!60, line width=0.8pt,
        minimum size=2cm, align=center, font=\small\bfseries, inner sep=3pt
    },
    helpful/.style={->, line width=0.7pt, black!70, bend left=18},
    selfloop/.style={->, line width=0.5pt, black!40, looseness=5},
    drift/.style={->, line width=0.5pt, black!30, dashed, bend left=15},
    actlbl/.style={font=\tiny, text=black!60, align=center},
    problbl/.style={font=\tiny\itshape, text=black!50},
]

% ── States ──
\node[state, fill=green!15] (flow)    at (0,0)      {FLOW};
\node[state, fill=yellow!20] (bore)   at (-4.5,-3.5) {BOREDOM};
\node[state, fill=orange!18] (anx)    at (4.5,-3.5)  {ANXIETY};
\node[state, fill=red!12]   (apa)     at (0,-5.5)    {APATHY};

% ── Helpful actions → FLOW (p=0.75) ──
\draw[helpful] (bore) to
    node[actlbl, above left, pos=0.45] {INCREASE\_URGENCY\\RAISE\_STAKES\\ADD\_MYSTERY}
    node[problbl, below left, pos=0.55] {$p=0.75$}
    (flow);
\draw[helpful] (anx) to
    node[actlbl, above right, pos=0.45] {PROVIDE\_GUIDANCE\\LOWER\_STAKES}
    node[problbl, below right, pos=0.55] {$p=0.75$}
    (flow);
\draw[helpful] (apa) to
    node[actlbl, right, pos=0.5] {ADD\_MYSTERY\\LORE\_REWARD}
    node[problbl, left, pos=0.5] {$p=0.75$}
    (flow);

% ── Self-loops (remain, p=0.25 or unhelpful) ──
\draw[selfloop] (bore) to[out=210, in=240] node[problbl, below] {$p=0.25$} (bore);
\draw[selfloop] (anx) to[out=-30, in=-60] node[problbl, below] {$p=0.25$} (anx);
\draw[selfloop] (apa) to[out=210, in=240] node[problbl, below] {$p=0.25$} (apa);

% ── FLOW self-loop (NO_CHANGE) ──
\draw[selfloop] (flow) to[out=60, in=90]
    node[actlbl, above] {NO\_CHANGE}
    node[problbl, right, xshift=3pt] {$p=0.90$}
    (flow);

% ── Drift from FLOW → BOREDOM ──
\draw[drift] (flow) to
    node[problbl, left, pos=0.5] {drift $p=0.10$}
    (bore);

\end{tikzpicture}
\caption{MDP state-transition diagram. Solid arrows show transitions caused by helpful narrative actions ($p = 0.75$ to FLOW). Self-loops indicate remaining in the current state. The dashed arrow shows the stochastic drift from FLOW to BOREDOM when no helpful action is needed.}
\label{fig:mdp_transitions}
\end{figure}

\textbf{Discount factor $\gamma$:} 0.95, reflecting moderate preference for near-term reward while maintaining long-horizon planning.

\subsection{Reward Function Design}

The reward function is defined over (prev\_state, new\_state, action, steps\_in\_state, last\_3\_actions) and comprises five components, detailed in Table~\ref{tab:reward}.

\begin{table}[H]
\centering
\caption{Reward function components and values.}
\label{tab:reward}
\small
\begin{tabular}{@{}lll@{}}
\toprule
\textbf{Component} & \textbf{Value} & \textbf{Trigger} \\
\midrule
State reward (FLOW) & $+1.0$ & new\_state == FLOW \\
State reward (BOREDOM) & $-0.3$ & new\_state == BOREDOM \\
State reward (ANXIETY) & $-0.5$ & new\_state == ANXIETY \\
State reward (APATHY) & $-0.8$ & new\_state == APATHY \\
Step penalty & $-0.05$ & Every step \\
Flow streak bonus & $+0.20$ & FLOW AND steps\_in\_state $>$ 2 \\
Transition bonus & $+0.30$ & prev $\in$ \{ANXIETY, BOREDOM\} AND new == FLOW \\
Repetition penalty & $-0.15$ & All 3 recent actions identical \\
\bottomrule
\end{tabular}
\end{table}

The reward function encodes a clear priority ordering. APATHY draws the heaviest penalty ($-0.8$) because it represents the deepest disengagement --- a player who has essentially given up. ANXIETY ($-0.5$) is penalised more than BOREDOM ($-0.3$) because an overwhelmed player is at higher risk of quitting the session entirely. The step penalty ($-0.05$) applies at every timestep regardless of state, creating constant pressure toward action efficiency and discouraging the agent from drifting into passive NO\_CHANGE loops. The Flow streak bonus ($+0.20$) rewards sustained optimal experience, not just momentary hits of Flow: it activates only after the player has been in FLOW for more than two consecutive steps, encouraging the agent to maintain conditions rather than simply reaching them. The transition bonus ($+0.30$) provides additional shaping reward for the key recovery transitions --- moving a player from ANXIETY or BOREDOM back into FLOW --- which are the most valuable moments for the system to get right. Finally, the repetition penalty ($-0.15$) fires when all three most recent actions are identical, discouraging degenerate policies in which the agent converges on a single favourite action and applies it regardless of context.

\subsection{Cold-Start Heuristic}

For the first 120~seconds of a session, or when the PPO model is unavailable, the \texttt{NarrativeAgent} falls back to a deterministic heuristic policy:

\begin{itemize}
\item ANXIETY $\rightarrow$ PROVIDE\_GUIDANCE
\item BOREDOM $\rightarrow$ INCREASE\_URGENCY
\item APATHY $\rightarrow$ ADD\_MYSTERY
\item FLOW $\rightarrow$ NO\_CHANGE
\end{itemize}

The 120-second threshold serves a dual purpose. It gives the player time to settle into the game world --- learning the controls, orienting to the environment, forming initial goals --- before the learned policy begins making adaptation decisions. And it gives the player modelling module time to accumulate enough telemetry for the Hexad profile and Flow classification to be meaningful. A PPO policy operating on a player model derived from ten seconds of data would be making decisions on noise; the heuristic ensures that the first two minutes of narrative adaptation are sensible even when the signal is thin.

\section{Unity Plugin Design}

The Unity plugin comprises three C\# components: \texttt{PlayerDataLogger}, \texttt{NarrativeManager}, and \texttt{ContentInjector}.

\texttt{PlayerDataLogger} is a MonoBehaviour that maintains running accumulators for all 23 telemetry fields. Every five seconds, it serializes the accumulated values into a \texttt{TelemetryBatch} JSON object and dispatches it to the FastAPI backend via an async HTTP POST. It exposes a public API for game systems to register events: \texttt{RecordKill()}, \texttt{RecordDeath()}, \texttt{RecordDialogueShown()}, and others. This event-based registration model decouples the telemetry system from specific game mechanics, making the plugin applicable to any Unity game that exposes these event hooks.

\texttt{NarrativeManager} polls the \texttt{/narrative/generate} endpoint following each telemetry transmission. \texttt{ContentInjector} maps \texttt{NarrativeContentType} to the appropriate UI mechanism: dialogue content is routed to the DialoguePanel, description content to the DescriptionPanel, and quest updates to the QuestLog. The injector implements a content queue with a 3-second minimum spacing between consecutive injections to prevent narrative overcrowding.

For the Control condition, the \texttt{NarrativeManager} is replaced by a \texttt{StaticNarrativeManager} that reads from a \texttt{StaticNarrativeContent} ScriptableObject and fires narrative events at scripted triggers with no player-state dependency.

\section{Testbed Game Design}

The testbed is a 3D RPG built in Unity 6 LTS using the AnyRPG Core framework \citep{anyrpg2024core}. AnyRPG provides a complete RPG engine --- combat, questing, inventory, NPC dialogue, and character progression --- out of the box, which allowed the development effort to focus on the narrative personalization system rather than on reimplementing standard RPG mechanics from scratch.

The game world is set in the Contested Vale, a once-fertile trade corridor now fought over by two rival factions. The setting was chosen to support a range of player behaviours: combat encounters for Achievers and Disruptors, exploration opportunities for Explorers and Free Spirits, NPC interactions for Socializers, and lore collectibles for Philanthropists. The world is described in structured lore documents (stored in the \texttt{lore/} directory) that double as the RAG corpus for the narrative generation pipeline.

Three navigable areas make up the game world:

\begin{enumerate}
\item \textbf{Commander Varis's Post} --- a fortified camp at the northern edge serving as the starting area. Varis is the primary quest giver: a methodical, unsentimental military commander who gives orders without explaining them. Players who engage extensively with Varis register high Achiever signals.

\item \textbf{The Contested Zone} --- a combat-focused area containing enemy units and environmental hazards. Players who clear this area efficiently register high Achiever and Disruptor signals; players who explore its periphery register Explorer signals.

\item \textbf{The Waystone / Chronicler's Camp} --- a neutral area where the Chronicler, a scholar who has observed the Vale for thirty years, provides lore and context. The presence of lore collectibles creates a natural Explorer/Philanthropist signal divergence.
\end{enumerate}

Three NPCs populate the world, each designed with a distinct voice and function that the LLM must maintain when generating dialogue for them. Commander Varis is the primary quest giver: a methodical, unsentimental military officer who gives orders and expects compliance. The Chronicler is the lore source: a scholar who has observed the Vale for thirty years and speaks in measured, academic prose. Sera the Merchant is the pragmatic intermediary: a supply trader whose dialogue is clipped, transactional, and occasionally sardonic. These three voices serve as a consistency test for the narrative generation pipeline --- if the LLM generates dialogue for the Chronicler that sounds like Varis, the player will notice.

The main quest (Clear the Advance) provides a linear spine through the combat zone, ensuring that all players encounter a minimum threshold of challenge and narrative content. A secondary collection quest (Salvage the Fallen) encourages off-path exploration and rewards lore engagement. A standard play session runs approximately 25--30 minutes, which is long enough for the player modelling module to build a meaningful profile but short enough to remain practical within the constraints of a controlled user study.

\section{Evaluation Design}

\subsection{Research Hypotheses}

The following directional hypotheses are tested:

\begin{description}
\item[H1:] Participants in the Experimental condition will score significantly higher on the GEQ Flow subscale than participants in the Control condition.
\item[H2:] Participants in the Experimental condition will score significantly higher on the GEQ Immersion subscale than participants in the Control condition.
\item[H3:] Participants in the Experimental condition will score significantly higher on the NPS-3 Narrative Personalization Scale than participants in the Control condition.
\end{description}

H1 is designated the primary hypothesis because Flow is the direct optimisation target of the adaptation engine's reward function. If the system works as designed, H1 is the hypothesis it is most likely to support; if the system fails, H1 is the hypothesis that will be most clearly refuted.

\subsection{Participants and Recruitment}

A total of $N = 50$ participants are recruited (25 Experimental, 25 Control) via convenience sampling through university mailing lists and gaming society channels. Inclusion criteria: aged 18 or over; self-reported experience playing RPG or action-adventure games; no prior exposure to the testbed game. An a priori power analysis for Welch's independent-samples $t$-test targeting $d = 0.5$ (medium effect), $\alpha = 0.05$ (two-tailed), and $1 - \beta = 0.80$ yields a required sample size of approximately 51 total participants \citep{faul2007gpower}, indicating that $N = 50$ is marginally below the strict threshold but remains reasonable given practical constraints.

Participants are randomly assigned to conditions using a balanced random allocation procedure. No block randomization by demographic is applied, given the expected homogeneity of the university convenience sample.

\subsection{Instruments}

\textbf{Game Experience Questionnaire (GEQ) --- Core Module.} The GEQ \citep{ijsselsteijn2013geq} is a 33-item instrument measuring seven dimensions of game experience on a 0--4 scale. The primary DV is the GEQ \textbf{Flow subscale} (items 5, 13, 25, 28, 31; items 5 and 13 reverse-scored). The secondary DV is the GEQ \textbf{Immersion subscale} (items 3, 12, 18, 19, 27, 30).

\textbf{miniPXI --- Player Experience Inventory (Short Form).} The miniPXI \citep{vornhagen2022minipxi} is an 11-item instrument on a 1--7 scale covering six player experience dimensions. The total score (mean of all 11 items) is used as a general player experience index.

\textbf{Narrative Personalization Scale (NPS-3).} The NPS-3 is a custom three-item scale on a 1--7 Likert scale: (1) ``The story felt tailored to how I was playing,'' (2) ``The narrative responded to my choices and actions,'' and (3) ``The game seemed to understand what kind of player I am.'' Cronbach's alpha will be reported; if $\alpha < 0.70$, items will be analysed individually.

\textbf{Behavioural Metrics (from telemetry).} Three behavioural metrics derived from telemetry serve as objective secondary measures: voluntary NPC interactions, dialogue read ratio, and exploration coverage.

\subsection{Procedure}

Each participant session follows a standardized procedure conducted individually:

\begin{enumerate}
\item \textbf{Arrival and consent} (5 min): briefing, participant information sheet, informed consent, audio recording consent.
\item \textbf{Tutorial} (5 min): condition-identical tutorial scene familiarizing participants with movement, combat, and NPC interaction.
\item \textbf{Game session} (25--30 min): participant plays in their assigned condition. The researcher is present but silent. Session is time-limited to 30 minutes.
\item \textbf{Questionnaires} (10 min): GEQ Core Module, miniPXI, and NPS-3 completed immediately post-session.
\item \textbf{Semi-structured interview} (20 min): following the interview guide (Appendix~\ref{app:interview}), audio recorded.
\item \textbf{Debrief} (5 min): explanation of the system and opportunity for questions.
\end{enumerate}

Total session duration: approximately 65 minutes per participant.

\subsection{Ethical Considerations}

This research received ethical approval through the institutional SPER process. All participants provide informed consent prior to participation and are informed of their right to withdraw at any time without penalty. Telemetry data is stored locally on the research machine in a SQLite database without personally identifiable information beyond the session identifier. Audio recordings from interviews are stored on an encrypted drive and transcribed with participant consent. All data is handled in accordance with GDPR requirements and university data protection policy. Participants are debriefed following the session and informed of the study's full purpose.

\subsection{Analysis Plan}

\textbf{Quantitative analysis} follows these steps:

\begin{enumerate}
\item \textbf{Normality assessment:} Shapiro-Wilk tests per DV per condition. If normality is violated ($p < 0.05$), Mann-Whitney U is substituted for Welch's $t$-test.
\item \textbf{Primary test (H1):} Welch's independent-samples $t$-test comparing GEQ Flow subscale scores. Effect size reported as Cohen's $d$ \citep{cohen1988statistical}.
\item \textbf{Secondary tests (H2, H3):} Identical procedure. Bonferroni correction applied across three hypotheses ($\alpha_{\text{adj}} = 0.0167$).
\item \textbf{Reliability:} Cronbach's alpha computed for each subscale.
\item \textbf{Behavioural metrics:} Welch's $t$-tests as exploratory analyses without correction.
\end{enumerate}

\textbf{Qualitative analysis} follows the six-phase Reflexive Thematic Analysis procedure \citep{braun2022thematic}: familiarization, coding, constructing themes, reviewing, defining themes, and writing up with representative quotes.
